\section{Mathematical Foundations}

\subsection{Linear Algebra}

\subsubsection{Matrix Operations and Decompositions}

Consider a system of linear equations represented in matrix form $Ax = b$, where
$A \in \mathbb{R}^{m \times n}$. The solution depends on the rank of~$A$.

The general least-squares solution uses the \textbf{pseudo-inverse}:

$$\begin{align}
x^* &= A^+ b \\
    &= (A^T A)^{-1} A^T b \\
    &= V \Sigma^+ U^T b
\end{align}$$

A piecewise function with the cases environment:

$$f(x) = \begin{cases}
x^2 & \text{if } x \geq 0 \\
-x^2 & \text{if } x < 0 \\
0 & \text{otherwise}
\end{cases}$$

The identity matrix and a general matrix side by side:

$$I_3 = \begin{matrix}
1 & 0 & 0 \\
0 & 1 & 0 \\
0 & 0 & 1
\end{matrix}
\quad
A = \begin{pmatrix}
a_{11} & a_{12} & a_{13} \\
a_{21} & a_{22} & a_{23} \\
a_{31} & a_{32} & a_{33}
\end{pmatrix}$$

\subsection{Probability Theory}

\subsubsection{Bayes' Theorem and Applications}

Given events $A$ and $B$ with $P(B) > 0$:

$$P(A \mid B) = \frac{P(B \mid A) \cdot P(A)}{P(B)}$$

The total probability expansion:

$$P(B) = \sum_{i=1}^{n} P(B \mid A_i) P(A_i)$$


\section{Programming Languages}

\subsection{Algorithm Implementations}

\subsubsection{Dynamic Programming Examples}

The classic coin change problem can be solved with DP. Below is the recurrence:

$$dp[i] = \begin{cases}
0 & \text{if } i = 0 \\
\min_{c \in \text{coins}} (dp[i - c] + 1) & \text{if } i > 0 \\
\infty & \text{if no valid coin}
\end{cases}$$

\begin{lstlisting}[language=Python]
def coin_change(coins: list[int], amount: int) -> int:
    dp = [float('inf')] * (amount + 1)
    dp[0] = 0
    for i in range(1, amount + 1):
        for c in coins:
            if c <= i and dp[i - c] + 1 < dp[i]:
                dp[i] = dp[i - c] + 1
    return dp[amount] if dp[amount] != float('inf') else -1
\end{lstlisting}

\subsubsection{Graph Algorithms in C++}

\begin{lstlisting}[language=C++]
#include <vector>
#include <queue>
using namespace std;

vector<int> bfs(const vector<vector<int>>& adj, int src) {
    int n = adj.size();
    vector<int> dist(n, -1);
    queue<int> q;
    dist[src] = 0;
    q.push(src);
    while (!q.empty()) {
        int u = q.front(); q.pop();
        for (int v : adj[u]) {
            if (dist[v] == -1) {
                dist[v] = dist[u] + 1;
                q.push(v);
            }
        }
    }
    return dist;
}
\end{lstlisting}

\subsubsection{Data Structures in Java}

\begin{lstlisting}[language=Java]
import java.util.*;

public class SegmentTree {
    private int[] tree;
    private int n;

    public SegmentTree(int[] arr) {
        n = arr.length;
        tree = new int[4 * n];
        build(arr, 1, 0, n - 1);
    }

    private void build(int[] arr, int node, int lo, int hi) {
        if (lo == hi) {
            tree[node] = arr[lo];
            return;
        }
        int mid = (lo + hi) / 2;
        build(arr, 2 * node, lo, mid);
        build(arr, 2 * node + 1, mid + 1, hi);
        tree[node] = Math.min(tree[2 * node], tree[2 * node + 1]);
    }

    public int query(int node, int lo, int hi, int l, int r) {
        if (r < lo || hi < l) return Integer.MAX_VALUE;
        if (l <= lo && hi <= r) return tree[node];
        int mid = (lo + hi) / 2;
        return Math.min(
            query(2 * node, lo, mid, l, r),
            query(2 * node + 1, mid + 1, hi, l, r)
        );
    }
}
\end{lstlisting}


\section{Reference Tables and Typography}

\subsection{Algorithm Complexity Reference}

\subsubsection{Comprehensive Comparison Table}

\begin{tabular}{|l|c|c|c|r|}
\hline
\multicolumn{5}{|c|}{\textbf{Algorithm Complexity Summary}} \\
\hline
\textbf{Algorithm} & \textbf{Best} & \textbf{Average} & \textbf{Worst} & \textbf{Space} \\
\hline
Merge Sort & $O(n \log n)$ & $O(n \log n)$ & $O(n \log n)$ & $O(n)$ \\
\hline
Quick Sort & $O(n \log n)$ & $O(n \log n)$ & $O(n^2)$ & $O(\log n)$ \\
\hline
\multirow{2}{*}{Heap Sort} & \multicolumn{2}{c|}{$O(n \log n)$} & $O(n \log n)$ & $O(1)$ \\
\cline{2-5}
 & \multicolumn{4}{c|}{\textit{In-place, not stable}} \\
\hline
Counting Sort & $O(n+k)$ & $O(n+k)$ & $O(n+k)$ & $O(k)$ \\
\hline
\multirow{2}{*}{Radix Sort} & \multicolumn{3}{c|}{$O(d \cdot (n+k))$} & $O(n+k)$ \\
\cline{2-5}
 & \multicolumn{4}{c|}{\textit{$d$ = digits, $k$ = base}} \\
\hline
\end{tabular}

\subsection{Nested List Structures}

\subsubsection{Algorithm Design Techniques}

\begin{enumerate}
\item \textbf{Divide and Conquer}
  \begin{itemize}
  \item Split the problem into subproblems
    \begin{enumerate}
    \item Identify independent substructures
    \item Determine base cases
    \item Design the merge step
    \end{enumerate}
  \item Solve each subproblem recursively
    \begin{itemize}
    \item Apply memoization if subproblems overlap
    \item Consider iterative bottom-up approaches
      \begin{itemize}
      \item Tabulation avoids stack overflow
      \item Often has better cache performance
      \item Easier to parallelize
      \end{itemize}
    \end{itemize}
  \end{itemize}
\item \textbf{Greedy Algorithms}
  \begin{itemize}
  \item Make locally optimal choices
  \item Prove the greedy choice property
    \begin{enumerate}
    \item Show optimal substructure
    \item Prove exchange argument
    \item Verify by contradiction
    \end{enumerate}
  \end{itemize}
\item \textbf{Network Flow}
  \begin{itemize}
  \item Model as a directed graph with capacities
  \item Apply Ford--Fulkerson or Dinic's algorithm
  \end{itemize}
\end{enumerate}

\subsection{Typography and Formatting Showcase}

\subsubsection{Text Formatting Commands}

This paragraph demonstrates \textbf{bold text}, \textit{italic text},
\underline{underlined text}, \texttt{monospace text}, \sout{struck-through text},
and \textsc{small caps text}.

Nested formatting: \textbf{\textit{bold italic}}, \textit{\texttt{italic monospace}},
\textbf{\underline{bold underline}}, and \textbf{\textit{\underline{bold italic underline}}}.

\subsubsection{Typographic Details}

Em dash---used for parenthetical statements---like this one.
En dash: pages 12--34, years 2020--2025.

``Smart double quotes'' and `smart single quotes' are supported.

A non-breaking~space prevents line breaks between words. The cost is \$42.00,
which is 15\% of the budget. The expression $a \& b$ uses the bitwise AND operator.
See item \#7 in the list. The underscore character \_\ is used in identifiers.
Curly braces \{ and \} delimit groups in \LaTeX.

\subsubsection{Links and References}

Visit \href{https://codeforces.com}{Codeforces} for competitive programming contests.
The full URL is \url{https://codeforces.com/problemset}.
Documentation is at \href{https://cp-algorithms.com/algebra/sieve-of-eratosthenes.html}{CP~Algorithms}.

\epigraph{The art of programming is the art of organizing complexity.}{Edsger W. Dijkstra}

More inline math to stress-test: the Euler identity $e^{i\pi} + 1 = 0$ connects
five fundamental constants. The binomial coefficient $\binom{n}{k} = \frac{n!}{k!(n-k)!}$
counts combinations. The integral $\int_0^\infty e^{-x^2} dx = \frac{\sqrt{\pi}}{2}$
is the Gaussian integral.

$$\begin{align}
\nabla \cdot \mathbf{E} &= \frac{\rho}{\varepsilon_0} \\
\nabla \cdot \mathbf{B} &= 0 \\
\nabla \times \mathbf{E} &= -\frac{\partial \mathbf{B}}{\partial t} \\
\nabla \times \mathbf{B} &= \mu_0 \mathbf{J} + \mu_0 \varepsilon_0 \frac{\partial \mathbf{E}}{\partial t}
\end{align}$$


\begin{diagram}[id="complexity-graph", label="Algorithm relationships"]
\shape{G}{Graph}{nodes=["DFS","BFS","Dijkstra","Bellman-Ford","Floyd"], edges=[("DFS","BFS"),("BFS","Dijkstra"),("Dijkstra","Bellman-Ford"),("Bellman-Ford","Floyd")], directed=true, layout="hierarchical"}
\recolor{G.node[DFS]}{state=current}
\recolor{G.node[Floyd]}{state=done}
\end{diagram}
